\documentclass[11pt,a4paper]{article}

\usepackage[T1]{fontenc}
\usepackage[utf8]{inputenc}
\usepackage[margin=25mm]{geometry}
\usepackage{lmodern}
\usepackage{microtype}
\usepackage{amsmath,amssymb}
\usepackage{booktabs,tabularx,array}
\usepackage{xcolor}
\usepackage{tikz}
\usetikzlibrary{arrows.meta,positioning,calc,fit,backgrounds,decorations.pathreplacing}
\usepackage[hidelinks]{hyperref}

\definecolor{routeblue}{HTML}{245B85}
\definecolor{parcelgreen}{HTML}{237A57}
\definecolor{warmorange}{HTML}{A65D17}
\definecolor{mutedgray}{HTML}{59636E}
\tikzset{
  flow/.style={-{Latex[length=2mm]},thick,draw=mutedgray},
  optionalflow/.style={flow,dashed},
  process/.style={draw=routeblue,rounded corners=2pt,fill=routeblue!5,
    align=center,font=\small,inner sep=6pt},
  routebox/.style={draw=mutedgray,rounded corners=2pt,align=center,
    font=\small,inner sep=5pt},
  passenger/.style={draw=routeblue,fill=routeblue!9,rounded corners=2pt,
    align=center,font=\small,inner sep=5pt},
  parcel/.style={draw=parcelgreen,fill=parcelgreen!9,rounded corners=2pt,
    align=center,font=\small,inner sep=5pt},
  note/.style={font=\footnotesize,align=center,text=mutedgray}
}
\newcolumntype{Y}{>{\raggedright\arraybackslash}X}
\newcommand{\code}[1]{\texttt{\detokenize{#1}}}
\setlength{\emergencystretch}{2em}

\title{SoICT 2026: Integrated Passenger and Parcel Transportation\\
  \large Problem Specification and Search Methodology}
\date{}

\begin{document}
\maketitle

\begin{abstract}
The method searches for taxi routes and then chooses routes for all taxis.
Bounded depth-first search builds routes visit by visit. Randomized solution
construction finds more routes through removal, greedy insertion, and local
search. Branch-and-bound combines stored routes without assigning a request
to two taxis. When there are many taxis with many stored routes, it repeatedly solves
smaller groups of taxis. Request prices guide the search toward useful alternatives. The
methodology first explains this architecture, then follows a small example
through the main search components. Formulas and implementation details are
kept in a separate technical part.
\end{abstract}

\tableofcontents
\clearpage
\part{Problem setting}
\input{sections/problem-statement}

\clearpage
\part{Algorithm methodology}
\label{part:methodology}
\input{sections/methodology}

\clearpage
\appendix
\part{Technical specification and implementation details}
\label{part:technical}
\input{sections/technical-details}

\end{document}
